\subsection{Comparison with the continuous profile: a uniform inverse bound and conclusion}\label{sec:P06}
We use (TR1)--(TR9). The constants depend on the bounds in Section~\ref{sec:P02} and
on the continuous profile, but are independent of $n$.

\paragraph{The exact averaged Jacobian.}
Differentiating the root in (TR5) gives, for each square,
\[
 w(z)_{ij}=\frac{\partial\rho(z)_{ij}}{\partial(\mathcal Tz)_{ij}}
 =\frac{\rho(z)_{ij}(1-\rho(z)_{ij}/n)}{1-4\rho(z)_{ij}/n}.
\]
On the segment in (TR9), $a\le w(z)\le C_w:=2Q$. The Jacobian is
\[
 J(z)h=\operatorname{diag}(0_R,0_C,u_D(z),u_A(z))h
                         +\mathcal R(w(z)\mathcal T h).
                                                        \tag{TP1}
\]
Its average along $z(t)=(1-t)z^0+tz^*$ has exactly the same form:
\[
 \overline J=\int_0^1J(z(t))\,dt
 =\operatorname{diag}(0_R,0_C,\overline u_D,\overline u_A)
                              +\mathcal R(\overline w\mathcal T),
 \quad a\le\overline w\le C_w,\quad\zeta\le\overline u_D,\overline u_A\le1.
                                                        \tag{TP2}
\]
It is self-adjoint with respect to the inner products defined between (TR3) and (TR4), annihilates $v_0$, and
preserves $E_n$.

\paragraph{Coercivity in the chosen gauge.}
We prove the estimate directly for the operator (TP2), which differs from the matrix evaluated at the
maximizer. The argument is that of Section~\ref{sec:P03}.
If $h\in E_n$, the means of $h_R$ and $h_C$ are equal. Hence
\[
 \|h_R\circ R+h_C\circ C\|_\square^2
 =\|h_R\|_{2,R}^2+\|h_C\|_{2,C}^2+2\overline h_R\overline h_C
 \ge\|h_R\|_{2,R}^2+\|h_C\|_{2,C}^2.
\]
Each diagonal lifting has $L^2$ norm at most one, since
a diagonal contains at most $n$ squares. Writing
$h_R\circ R+h_C\circ C=\mathcal T h-h_D\circ D-h_A\circ A$
and using $(s-d-a)^2\le2s^2+4(d^2+a^2)$ gives
\[
 \|h\|_\ell^2\le2\|\mathcal T h\|_\square^2
                    +5(\|h_D\|_{2,D}^2+\|h_A\|_{2,A}^2).
\]
Since $\mathcal R^*=\mathcal T$, (TP2) implies
\[
 \langle h,\overline Jh\rangle_\ell
 \ge a\|\mathcal T h\|_\square^2
             +\zeta(\|h_D\|_{2,D}^2+\|h_A\|_{2,A}^2)
 \ge\lambda\|h\|_\ell^2,
 \quad\lambda=\min(a/2,\zeta/5)>0.
                                                        \tag{TP3}
\]
By finite dimensionality, $\overline J:E_n\to E_n$ is invertible and
$\|\overline J^{-1}k\|_\ell\le\lambda^{-1}\|k\|_\ell$.

\paragraph{From $L^2$ to the supremum norm by discrete transversality.}
Let $\overline Jh=k\in E_n$. The equation for a line $\ell$ in family
$f$ is
\[
 \left(d_{f,\ell}+\frac1n\sum_{ij\in\ell}\overline w_{ij}\right)h_{f,\ell}
 +\sum_{g\ne f}\frac1n\sum_{ij\in\ell}\overline w_{ij}h_{g,\ell_g(i,j)}
 =k_{f,\ell},
                                                        \tag{TP4}
\]
with $d_R=d_C=0$, $d_D=\overline u_D$, and $d_A=\overline u_A$.
The coefficient of the line's own component is at least $a_0:=\min(a,\zeta)>0$:
rows and columns contain exactly $n$ squares, and diagonals and antidiagonals satisfy
$d_{f,\ell}\ge\zeta$.
Two lines from distinct families meet in at most one square.
Thus, for $g\ne f$, the indices $\ell_g(i,j)$ encountered along $\ell$
are all distinct. Cauchy--Schwarz gives
\[
 \frac1n\sum_{ij\in\ell}|h_{g,\ell_g(i,j)}|
 \le\sqrt{\frac{|\ell|}{n}}
       \left(\frac1n\sum_{\ell'\text{ visited}}|h_{g,\ell'}|^2\right)^{1/2}
 \le\|h_g\|_{2,g}.
                                                        \tag{TP5}
\]
In the coupling between diagonals and antidiagonals, only indices of one parity are encountered,
but they are still distinct, so the same bound holds without any additional factor.
By (TP4)--(TP5), summing over the three transverse families yields
\[
 \|h\|_\infty\le a_0^{-1}
       (\|k\|_\infty+\sqrt3 C_w\|h\|_\ell).
\]
Since there are $6n-2$ lines, $\|k\|_\ell\le\sqrt6\|k\|_\infty$.
Together with (TP3), this gives the required uniform inverse bound:
\[
 \|\overline J^{-1}k\|_\infty\le K\|k\|_\infty,
 \qquad K=a_0^{-1}(1+\sqrt{18}C_w/\lambda).
                                                        \tag{TP6}
\]
The estimate holds for arbitrary data in $E_n$, including oscillatory data.

\begin{proposicion}
For all sufficiently large $n$, the canonical maximizer satisfies
\[
 \max_{ij}|n\pi_{ij}-\nu_1(x_i,x_j)|\le C/n,
 \quad C=4C_wKC_0+12C_\nu^2,
                                                        \tag{TP7}
\]
\[
 \max_k|h^*_{D,k}-u_D(k/n)|\le KC_0/n,\qquad
 \max_l|h^*_{A,l}-u_A(l/n)|\le KC_0/n.
                                                        \tag{TP8}
\]
\end{proposicion}
\begin{proof}
By the fundamental theorem of calculus and (TR5), we have the exact identity
$-F_n(z^0)=F_n(z^*)-F_n(z^0)=\overline J(z^*-z^0)$.
Both potentials belong to $E_n$. By (TR8) and (TP6),
$\|z^*-z^0\|_\infty\le KC_0/n$.
Along the segment, $d\rho/d(\mathcal Tz)=w\le C_w$ and
$\|\mathcal T\|_{\infty\to\infty}\le4$. Hence
$\|\rho(z^*)-\rho(z^0)\|_\infty\le4C_wKC_0/n$.
Now (TR7) and the identity $\rho(z^*)=n\pi$ prove (TP7).
The exponential vacancies are at most one along the segment,
so the scalar mean value theorem gives (TP8).
\end{proof}

Thus (TP7) establishes (LC3) with the same midpoint sampling.
